\documentclass{sysupptthesis}

\AddToHook{shipout/foreground}{\BackgroundPicture}
\author{Wang Kaiqi}
\title{Study of Radioactive Source Monitoring\\Based on ElasticNet Regression Method}
\subtitle{Bachelor's Thesis Defense Report}
\institute{Sino-French Institute of Nuclear Engineering and Technology, Sun Yat-sen University}
\date{April 2026}

\begin{document}

% ============ Title Page ============
\begin{frame}
    \titlepage
    \vspace*{-0.5cm}
    \begin{center}
        \Large\textbf{Sun Yat-sen University}\\[0.3em]
        \normalsize Supervisors: Prof.\ Shi Wei (CityU HK) \& Prof.\ Wang Kai (SYSU)
    \end{center}
\end{frame}

% ============ Outline ============
\begin{frame}
    \tableofcontents[sectionstyle=show,subsectionstyle=show/shaded/hide,subsubsectionstyle=show/shaded/hide]
\end{frame}

% ============================================================
\section{Research Background \& Motivation}
% ============================================================

\begin{frame}{Research Background}
    \begin{block}{Nuclear Facility Decommissioning}
        \begin{itemize}
            \item Global decommissioning projects increasing (Fukushima Daiichi, etc.)
            \item IAEA: Accurate source distribution is prerequisite for optimized decommissioning
            \item Ensuring personnel safety and minimizing waste generation
        \end{itemize}
    \end{block}

    \begin{block}{Three-Dimensional Complexity in Confined Spaces}
        \begin{itemize}
            \item \textbf{Composite source types}: localized hot-spots + continuous activation zones + dispersed media
            \item \textbf{Nuclide diversity}: $^{134}$Cs, $^{137}$Cs, $^{60}$Co, $^{131}$I with distinct energy spectra
            \item \textbf{Measurement constraints}: limited detectors, limited time, dose exposure limits
        \end{itemize}
    \end{block}
\end{frame}

\begin{frame}{Problem Statement}
    \begin{block}{The Inverse Problem}
        \begin{itemize}
            \item \textbf{Forward}: Given sources $\mathbf{s}$, predict detector readings $\mathbf{d}$
            \item \textbf{Inverse}: Given readings $\mathbf{d}$, reconstruct sources $\mathbf{s}$
        \end{itemize}
    \end{block}

    \defination{Mathematical Formulation}{
        \vspace{0.3em}
        \[
            \mathbf{d} = A \cdot \mathbf{s} + \boldsymbol{\varepsilon}
        \]
        where $\mathbf{d}$ = measurements, $A$ = contribution matrix, $\mathbf{s}$ = source vector, $\boldsymbol{\varepsilon}$ = noise.
    }

    \begin{block}{Key Challenges}
        \begin{itemize}
            \item \textbf{Underdetermined}: fewer detectors than unknowns ($m \ll n$)
            \item \textbf{Ill-posed}: noisy data leads to unstable solutions
            \item \textbf{Mixed}: sparse hot spots + continuous activation regions
        \end{itemize}
    \end{block}
\end{frame}

% ============================================================
\section{Literature Review}
% ============================================================

\begin{frame}{Existing Methods and Limitations}
    \begin{table}
        \centering
        \small
        \begin{tabular}{cccccl}
            \hline
            \textbf{Method} & \textbf{Sparse} & \textbf{Continuous} & \textbf{Multi-nuclide} & \textbf{Speed} & \textbf{Core Limitation} \\
            \hline
            ML-EM/MPMI & Fair & Good & Poor & Slow & $O(n^3)$; noise sensitive \\
            IRT-DEGL & Fair & Fair & Fair & V.Slow & Global optimization cost \\
            LASSO & \textbf{Excellent} & Poor & Good & Fast & Over-compresses continuous \\
            Fused LASSO & Fair & Good & Good & Medium & Needs neighborhood graph \\
            \textbf{ElasticNet} & \textbf{Excellent} & \textbf{Excellent} & Good & Fast & Not yet applied to nuclear \\
            \hline
        \end{tabular}
        \caption{Comparison of radioactive source reconstruction methods}
    \end{table}
\end{frame}

\begin{frame}{Why ElasticNet?}
    \begin{block}{Advantages over LASSO and Fused LASSO}
        \begin{itemize}
            \item LASSO: discrete/sparse distribution $\rightarrow$ precision loss for continuous regions
            \item Fused LASSO: requires predefined neighborhood graph + ADMM convergence issues
            \item ElasticNet: L1+L2 synergy without extra spatial constraints
        \end{itemize}
    \end{block}

    \begin{block}{Cross-Domain Evidence}
        \begin{itemize}
            \item Nuclear medicine (SPECT): outperforms pure L1, L2, fixed-weight EN
            \item XLCT imaging: superior localization + dual-source resolution
            \item Fluorescence molecular tomography: 30\%+ error reduction
        \end{itemize}
    \end{block}

    \theorem{Research Gap}{
        ElasticNet has \textbf{not been applied} to nuclear radioactive source monitoring --- this thesis fills that gap.
    }
\end{frame}

% ============================================================
\section{Theoretical Basis}
% ============================================================

\begin{frame}{ElasticNet Regression Formulation}
    \defination{Objective Function (Zou \& Hastie, 2005)}{
        \vspace{0.3em}
        \[
            \min_{\mathbf{w}} \quad \| \mathbf{y} - X\mathbf{w} \|_2^2 + \alpha \left[ \rho \| \mathbf{w} \|_1 + \frac{1-\rho}{2} \| \mathbf{w} \|_2^2 \right]
        \]
    }

    \begin{block}{L1 vs.\ L2 Penalty Roles}
        \begin{itemize}
            \item \textbf{L1} ($\rho \| \mathbf{w} \|_1$): promotes sparsity, drives irrelevant coefficients to zero $\rightarrow$ hot-spot localization
            \item \textbf{L2} ($\frac{1-\rho}{2} \| \mathbf{w} \|_2^2$): groups correlated variables, suppresses oscillation $\rightarrow$ continuous region reconstruction
            \item $\rho \to 1$: degrades to LASSO; \quad $\rho \to 0$: degrades to Ridge
        \end{itemize}
    \end{block}

    \begin{block}{Key Properties}
        \begin{itemize}
            \item Strictly convex when $\lambda_2 > 0$ $\rightarrow$ unique global minimum guaranteed
            \item No ADMM convergence issues (unlike Fused LASSO)
            \item Solved by Coordinate Descent algorithm
        \end{itemize}
    \end{block}
\end{frame}

\begin{frame}{Physical Interpretation}
    \begin{block}{L1 Penalty Physical Meaning}
        \begin{itemize}
            \item Count-type penalty on ``how many grid cells contain sources''
            \item Solution tends toward few non-zero concentrated elements --- sparse hot-spot localization
        \end{itemize}
    \end{block}

    \begin{block}{L2 Penalty Physical Meaning}
        \begin{itemize}
            \item Suppresses individual coefficient amplitude, promotes neighboring cells to have similar values
            \item Inside continuous activation regions: naturally smooth reconstruction (not ``discrete spots'' like LASSO)
        \end{itemize}
    \end{block}

    \theorem{Synergy Effect}{
        \begin{itemize}
            \item Near sparse hotspots: L1 dominates $\rightarrow$ precise localization
            \item Inside continuous regions: L2 dominates $\rightarrow$ smooth morphology
            \item Transition zones: both cooperate $\rightarrow$ natural morphological transition
        \end{itemize}
    }
\end{frame}

% ============================================================
\section{Algorithm Implementation \& Optimization}
% ============================================================

\begin{frame}{Algorithm Pipeline}
    \begin{block}{Stage 1: Preprocessing}
        \begin{itemize}
            \item Column normalization: scale each column of $A$ to unit $L_2$ norm
            \item Non-negative SVD matrix construction
            \item Condition number assessment for adaptive $\alpha$ range
        \end{itemize}
    \end{block}

    \begin{block}{Stage 2: ElasticNet Optimization}
        \begin{itemize}
            \item \texttt{ElasticNet(positive=True)} --- hard non-negativity constraint
            \item GridSearchCV: $\alpha \in [10^{-7}, 10^{0}]$ (50 values), $\rho \in [0.01, 0.99]$ (80 values)
            \item 5-fold cross-validation with negative MSE scoring
            \item Custom intercept penalty for $|\beta_0 / \bar{y}| > 10\%$
        \end{itemize}
    \end{block}

    \begin{block}{Stage 3: Post-processing}
        \begin{itemize}
            \item Denormalize: $\mathbf{w}_{\text{est}} = \mathbf{w}_{\text{norm}} \oslash \text{col\_norms}$
            \item Non-negative truncation, intercept analysis
        \end{itemize}
    \end{block}
\end{frame}

\begin{frame}{Key Optimization Strategies}
    \begin{block}{Adaptive $\alpha$ Search Range}
        \begin{itemize}
            \item cond $= 100$: $\alpha \in [10^{-6}, 10^{0}]$
            \item cond $= 500$: $\alpha \in [10^{-5}, 10^{1}]$
            \item cond $= 1000$: $\alpha \in [10^{-4}, 10^{2}]$
            \item Optimal \texttt{l1\_ratio} $\approx 0.99$ $\rightarrow$ strong sparsity prior confirmed
        \end{itemize}
    \end{block}

    \begin{block}{Non-negative Constraint Enhancement}
        \begin{itemize}
            \item \texttt{positive=True}: hard non-negative constraint
            \item Non-negative SVD matrix construction eliminates negative response coefficients
            \item ``Non-negativity + Sparse Selection + Smooth Anti-noise'' triple mechanism
        \end{itemize}
    \end{block}

    \begin{block}{Intercept Fitting Analysis}
        \begin{itemize}
            \item $|\beta_0/\bar{y}| < 10\%$: improves generalization
            \item $|\beta_0/\bar{y}| \gg 10\%$: reconstruction degrades
            \item ElasticNet has an \textbf{upper limit of noise tolerance}
        \end{itemize}
    \end{block}
\end{frame}

% ============================================================
\section{PHITS Monte Carlo Simulation}
% ============================================================

\begin{frame}{PHITS Simulation Configuration}
    \begin{block}{PHITS (JAEA): Monte Carlo Particle Transport}
        \begin{itemize}
            \item $10^8$ particles $\times$ 100 batches for statistical precision
            \item Cs-137 point sources, 0.662 MeV $\gamma$-ray
        \end{itemize}
    \end{block}

    \begin{table}
        \centering
        \small
        \begin{tabular}{ll}
            \hline
            \textbf{Parameter} & \textbf{Value} \\
            \hline
            Simulation space & Spherical (R=10m, reactor-centered) \\
            Radioactive sources & 10 Cs-137 point sources, 1 Ci each \\
            Detectors & 8 spherical detectors (R=25mm) at $z=20$mm \\
            Grid division & $4 \times 3 = 12$ cells \\
            Medium & Air (N, O, Ar at standard density) \\
            Contribution matrix & $8 \times 12$ (only 11/96 non-zero) \\
            Matrix rank & 6/8 \\
            \hline
        \end{tabular}
        \caption{PHITS simulation parameters}
    \end{table}
\end{frame}

% ============================================================
\section{Results \& Analysis}
% ============================================================

\begin{frame}{PHITS Validation Results}
    \begin{block}{Best Parameters: $\alpha = 1.0 \times 10^{-6}$, \texttt{l1\_ratio} $= 0.99$, $R^2 = 0.463$}
    \end{block}

    \begin{table}
        \centering
        \small
        \begin{tabular}{cccc}
            \hline
            \textbf{Cell} & \textbf{True Activity (GBq)} & \textbf{Estimate (GBq)} & \textbf{Rel.\ Error} \\
            \hline
            Cell 1  & $1.265 \times 10^{5}$ & $1.265 \times 10^{5}$ & \textbf{0.0\%} \\
            Cell 5  & $5.266 \times 10^{4}$ & $3.93 \times 10^{2}$ & \textcolor{red}{99.3\%} \\
            Cell 6  & $1.192 \times 10^{5}$ & $2.31 \times 10^{3}$ & \textcolor{red}{98.1\%} \\
            Cell 7  & $5.894 \times 10^{4}$ & $5.894 \times 10^{4}$ & \textbf{0.0\%} \\
            Cell 8  & $5.724 \times 10^{4}$ & $5.724 \times 10^{4}$ & \textbf{0.0\%} \\
            Cell 9  & $4.870 \times 10^{4}$ & $4.870 \times 10^{4}$ & \textbf{0.0\%} \\
            Cell 10 & $1.352 \times 10^{5}$ & $1.373 \times 10^{5}$ & \textbf{1.2\%} \\
            \hline
        \end{tabular}
        \caption{7 non-zero cells detected, 5 with $<2\%$ error}
    \end{table}
\end{frame}

\begin{frame}{Error Source Analysis}
    \begin{block}{Cell 5 \& Cell 6 Failure --- Physical, Not Algorithmic}
        \begin{itemize}
            \item Cell 6: only Detector 1 responds, extremely weak signal ($10^{-9}$ level)
            \item Cell 5: only Detector 5 responds, weak signal ($10^{-10}$ level)
            \item Both overwhelmed by stronger responses from other cells
        \end{itemize}
    \end{block}

    \theorem{Matrix Rank Analysis}{
        \begin{itemize}
            \item Numerical rank of $A$ = 6/8
            \item Only 6 of 12 cells algebraically distinguishable
            \item Physical indistinguishability from detector geometry, \textbf{not algorithm failure}
        \end{itemize}
    }

    \begin{block}{Implication}
        Improved detector layout would directly resolve these errors. ElasticNet is NOT the bottleneck.
    \end{block}
\end{frame}

\begin{frame}{Condition Number \& Bayesian Optimization}
    \begin{block}{Difficulty Inflection Points at Cond $\approx$ 150 and 300}
        \begin{itemize}
            \item Not simply ``more ill-conditioned = harder'' --- critical spectral structure
            \item All optimization strategies show diminishing returns at high condition numbers
            \item Optimal \texttt{l1\_ratio} $\approx 0.99$ universally --- confirms strong sparsity prior
        \end{itemize}
    \end{block}

    \begin{table}
        \centering
        \small
        \begin{tabular}{lcc}
            \hline
            & \textbf{Grid Search} & \textbf{Bayesian Opt.} \\
            \hline
            Evaluations & 150--450 & 40 \\
            Computation time & Baseline & $5\text{--}10\times$ faster \\
            Accuracy (low cond) & Reference & Comparable \\
            Accuracy (high cond) & \textbf{Higher ceiling} & Slightly worse \\
            \hline
        \end{tabular}
        \caption{Grid search vs.\ Bayesian optimization --- complementary, not substitutive}
    \end{table}
\end{frame}

% ============================================================
\section{Conclusions \& Future Work}
% ============================================================

\begin{frame}{Conclusions}
    \begin{block}{1. ElasticNet fills the application gap}
        Effectively handles ``sparse + continuous'' composite source reconstruction in nuclear monitoring.
    \end{block}

    \begin{block}{2. Optimization strategies improve accuracy significantly}
        Column normalization + adaptive $\alpha$ + non-negative constraint + intercept penalty.
    \end{block}

    \begin{block}{3. PHITS validation confirms practical value}
        All 7 non-zero cells detected; 5 cells with $<2\%$ error.
    \end{block}

    \begin{block}{4. Bayesian optimization: complementary to grid search}
        70--90\% computation savings; grid search has higher accuracy ceiling.
    \end{block}
\end{frame}

\begin{frame}{Future Work}
    \begin{block}{Mobile Detection Platforms (UGV / UAV)}
        \begin{itemize}
            \item Carry detectors for large-scale autonomous patrol
            \item Richer contribution matrix $\rightarrow$ improved reconstruction stability
        \end{itemize}
    \end{block}

    \begin{block}{Algorithm Improvements}
        \begin{itemize}
            \item Adaptive ElasticNet for automatic variable selection
            \item Bayesian ElasticNet with uncertainty quantification
            \item Deep learning + ElasticNet model cascade
        \end{itemize}
    \end{block}

    \begin{block}{Extended Applications}
        \begin{itemize}
            \item Multi-energy source discrimination
            \item 3D volumetric source reconstruction
            \item Real-time online reconstruction system
        \end{itemize}
    \end{block}
\end{frame}

% ============================================================
\section{References}
% ============================================================

\begin{frame}[allowframebreaks]{References}
    \small
    \begin{thebibliography}{99}
        \bibitem{tibshirani1996} Tibshirani R. Regression shrinkage and selection via the lasso. \textit{JRSS-B}, 1996.
        \bibitem{hoerl1970} Hoerl \& Kennard. Ridge regression. \textit{Technometrics}, 1970.
        \bibitem{sugaya2019} Sugaya et al. Inverse estimation for fuel debris search. \textit{ANE}, 2019.
        \bibitem{shi2023a} Shi W et al. LASSO with spectral information. \textit{ANE}, 2023.
        \bibitem{shi2023b} Shi W et al. LASSO: Theory and demonstration. \textit{PNE}, 2023.
        \bibitem{yamada2024} Yamada et al. Fused LASSO for radioactive materials. \textit{WM2024}, 2024.
        \bibitem{zou2005} Zou \& Hastie. Regularization via the elastic net. \textit{JRSS-B}, 2005.
        \bibitem{kasai2025} Kasai \& Otsuka. Dynamic elastic net for nuclear medicine. \textit{J. Imaging}, 2025.
        \bibitem{zhao2021} Zhao et al. Robust elastic net for XLCT. \textit{Phys. Med. Biol.}, 2021.
        \bibitem{sato2018} Sato et al. PHITS version 3.02. \textit{JNST}, 2018.
    \end{thebibliography}
\end{frame}

% ============================================================
\section{Q \& A}
% ============================================================

\begin{frame}
    \begin{center}
        {\Huge\textbf{Questions \& Discussion}}\\[2em]
        {\Large Thank you for your attention!}\\[1em]
        {\normalsize Wang Kaiqi | 22355080}\\[0.3em]
        {\small Supervisors: Prof.\ Shi Wei (CityU HK) \& Prof.\ Wang Kai (SYSU)}
    \end{center}
\end{frame}

\end{document}
